\chapter{Fundamentação teórica}
Este capítulo estabelece as bases conceituais necessárias para a compreensão do escalonamento adaptativo de compactações em árvores LSM sobre Memória Persistente. Inicialmente, discutem-se os fundamentos das estruturas Log-Structured Merge-trees, seguidos pela análise detalhada da mecânica de compactação e suas implicações de custo. Posteriormente, abordam-se as características da Memória Persistente (PMem) e como ela altera os paradigmas de armazenamento tradicionais. Por fim, revisam-se os trabalhos relacionados e as lacunas na literatura que motivam a abordagem adaptativa proposta.

\section{Bancos de Dados Key-Value e o Problema de Armazenamento Eficiente}

Sistemas de gerenciamento de dados modernos enfrentam um desafio fundamental, como
persistir e recuperar grandes volumes de informação com alta velocidade, baixa
latência e custo operacional razoável. Enquanto bancos de dados relacionais
tradicionais organizam dados em tabelas com esquemas rígidos e utilizam estruturas
como B-Trees para indexação em disco, esse modelo apresenta limitações sérias em
cenários de alta taxa de gravação, onde o acesso aleatório ao disco imposto pela
B-Tree para operações de \textit{insert} e \textit{update} torna-se um gargalo
à medida que o volume de escritas cresce~\cite{database_internals}.

Para contornar esse problema, sistemas orientados a pares chave-valor
(\textit{key-value stores}) como o RocksDB~\cite{rocksdb} e o
LevelDB~\cite{leveldb} adotam uma abordagem radicalmente
diferente: \textbf{nunca modificar dados diretamente em disco}. Em vez disso,
todas as operações, inclusive atualizações e deleções, são tratadas como
novas gravações sequenciais, o que maximiza a utilização da largura de banda do
dispositivo de armazenamento e elimina o I/O aleatório custoso~\cite{lsm_survey}.

Esse paradigma, chamado de \textit{append-only}, é a base da estrutura
\textit{Log-Structured Merge-Tree} (LSM-tree), descrita em detalhes na
Seção~\ref{sec:lsmtree}. A LSM-tree é hoje o padrão de fato na camada de
armazenamento de sistemas distribuídos de larga escala, sendo empregada em
plataformas de streaming, bancos de dados de séries temporais e motores de busca,
entre outros~\cite{rocksdb,lsm_survey}. Sua adoção em ambientes que utilizam
Memória Persistente, abordada na Seção~\ref{sec:pmem}, abre novas
oportunidades de otimização, mas também expõe fragilidades que este trabalho
busca otimizar.

\section{LSM-Trees}
\label{sec:lsmtree}

A estrutura \textit{Log-Structured Merge-tree} (LSM-tree), proposta originalmente
por~\citeauthor{ONeil1996TheLM}~\cite{ONeil1996TheLM}, tornou-se o padrão de fato
para a camada de armazenamento em sistemas NoSQL modernos e \textit{key-value
stores}, priorizando o desempenho de escrita por meio de operações
\textit{append-only}, evitando atualizações \textit{in-place} que geram I/O
aleatório custoso em disco~\cite{database_internals,lsm_survey}.

\subsection*{Arquitetura e Fluxo de Escrita}

A arquitetura de uma LSM-tree opera em múltiplas camadas, divididas entre memória
volátil e armazenamento persistente. Os principais componentes são~\cite{database_internals,rocksdb}:

\begin{itemize}
    \item \textbf{MemTable:} Um buffer em memória (geralmente uma SkipList ou
    árvore balanceada) que recebe todas as escritas. Como reside em memória
    volátil, as operações são rápidas. Para garantir durabilidade, as escritas são
    registradas sequencialmente em um \textit{Write-Ahead Log} (WAL) antes de
    entrarem na MemTable.

    \item \textbf{SSTables (\textit{Sorted String Tables}):} Quando a MemTable
    atinge um limite de tamanho, ela é convertida em um formato imutável, ordenado
    por chave, e despejada (\textit{flushed}) para o armazenamento persistente.
    Uma vez escrito, esse arquivo nunca é modificado.

    \item \textbf{Níveis (\textit{Levels}):} As SSTables são organizadas em
    níveis (L0, L1, \ldots, Ln). O nível L0 contém os arquivos recém-criados pelo
    \textit{flush}. À medida que os dados envelhecem e são compactados, eles
    migram para níveis inferiores, exponencialmente maiores.

    \item \textbf{Bloom Filters:} Para evitar leituras desnecessárias em disco,
    cada SSTable geralmente possui um Bloom Filter associado, que permite verificar
    rapidamente se uma chave definitivamente \emph{não} existe em um arquivo,
    economizando I/O de leitura.
\end{itemize}

\subsection*{O Trade-off Central: Escrita vs.\ Leitura}

A imutabilidade das SSTables é o que garante o alto desempenho de escrita, mas
introduz uma complexidade inerente nas leituras: uma mesma chave pode ter múltiplas
versões distribuídas por diferentes SSTables em diferentes níveis. Para ler um
valor, o sistema precisa consultar os níveis do mais recente (L0) ao mais antigo,
até encontrar a versão válida. Esse custo é mitigado pelos Bloom Filters e pela
garantia de ordem dentro de cada nível, mas representa um trade-off fundamental
no design da estrutura~\cite{Sarkar_2021,lsm_survey}.

O processo de \textbf{compactação} é o mecanismo que reequilibra esse trade-off:
ele consolida SSTables de diferentes níveis, descartando versões obsoletas e
reorganizando os dados para reduzir o custo de leitura futura. A mecânica
detalhada, as políticas de compactação e as métricas de custo associadas são
analisadas na Seção~\ref{sec:compaction}.

\subsection*{Compactação}\label{sec:compaction}

A compactação (\textit{compaction}) é o mecanismo fundamental que permite às
LSM-trees manterem desempenho sustentável ao longo do tempo. Como as SSTables são
imutáveis e novas versões de chaves são continuamente inseridas em níveis mais
recentes, o sistema naturalmente acumula múltiplas versões do mesmo dado,
além de registros de deleção (\textit{tombstones}). Sem intervenção, esse acúmulo
aumentaria progressivamente o custo de leitura e o consumo de espaço em disco.

De forma geral, a compactação consiste em selecionar um conjunto de SSTables,
tipicamente de um nível $L_i$ e possivelmente do nível subsequente $L_{i+1}$,
mesclá-las (\textit{merge}) de maneira ordenada e produzir um novo conjunto de
SSTables livres de redundâncias. Durante esse processo, versões obsoletas de
chaves são descartadas, e \textit{tombstones} podem ser eliminados quando já não
há versões antigas a serem invalidadas~\cite{database_internals,lsm_survey}.

O processo segue essencialmente um padrão de \textit{merge sort}: dado que cada
SSTable já está ordenada por chave, a fusão pode ser realizada de forma sequencial,
maximizando a localidade de acesso e evitando I/O aleatório. O resultado é um
novo conjunto de arquivos maiores, mais organizados e com menor sobreposição de
intervalos de chaves, o que reduz o número de estruturas que precisam ser
consultadas durante leituras futuras.

Existem diferentes políticas de compactação, sendo as mais comuns:

\begin{itemize}
    \item \textbf{Size-Tiered Compaction (STC):} Agrupa SSTables de tamanhos
    semelhantes dentro de um mesmo nível e as mescla quando um determinado número
    de arquivos é atingido. Essa abordagem favorece a redução do número de
    arquivos, mas pode manter maior sobreposição de chaves, impactando leituras.

    \item \textbf{Leveled Compaction (LC):} Organiza os dados em níveis com
    limites de tamanho bem definidos e mínima sobreposição entre SSTables dentro
    de um mesmo nível. A cada compactação, dados são promovidos de $L_i$ para
    $L_{i+1}$ de forma controlada. Essa política reduz significativamente o custo
    de leitura, ao custo de maior volume de escrita~\cite{rocksdb,lsm_survey}. É o padrão em sistemas como o RocksDB~\cite{rocksdb} e LevelDB~\cite{leveldb}.
\end{itemize}

Embora essencial, a compactação introduz custos relevantes, frequentemente
descritos em termos de três métricas principais:

\begin{itemize}
    \item \textbf{Amplificação de Escrita (Write Amplification - WAF):} Refere-se ao fato de que um mesmo dado pode ser reescrito múltiplas vezes ao longo de sua vida útil, à medida que é
    propagado entre níveis durante compactações.

    \item \textbf{Amplificação de Leitura (Read Amplification - RA):} Relaciona-se ao número de SSTables que precisam ser consultadas para responder a uma leitura. A compactação reduz
    esse custo ao consolidar dados.

    \item \textbf{Amplificação de Espaço (Space Amplification - SA):} Decorre da presença temporária de dados duplicados ou obsoletos antes que sejam removidos por compactação.
\end{itemize}

Essas três métricas formam um trade-off intrínseco: políticas que minimizam amplificação de escrita, como \textit{leveled compaction}, tendem a aumentar
\textit{write amplification}, enquanto abordagens mais simples, como
\textit{size-tiered}, apresentam comportamento oposto. A escolha da estratégia de
compactação, bem como seus parâmetros (tamanho de níveis, gatilhos e paralelismo),
tem impacto direto no desempenho e na eficiência do sistema~\cite{Sarkar_2021}.

Além disso, a compactação é uma das principais fontes de carga de I/O em sistemas
baseados em LSM-tree, podendo competir com operações de leitura e escrita
\textit{foreground}. Por esse motivo, implementações modernas, como RocksDB, utilizam
mecanismos de agendamento, limitação de taxa (\textit{rate limiting}) e execução
em segundo plano para mitigar interferências e manter previsibilidade de
desempenho~\cite{rocksdb}.

\subsection*{LSM-Trees sobre Memória Persistente}

As características operacionais das LSM-trees, em especial a alta frequência de
gravações físicas induzida pelo processo de compactação, ganham uma nova dimensão
quando o substrato de armazenamento é uma Memória Persistente
(PMem)~\cite{revisting_design_lsm,FlatLSM,SLM-DB}. A PMem, descrita em detalhes
na Seção~\ref{sec:pmem}, apresenta latências significativamente menores que SSDs,
mas impõe restrições de largura de banda de escrita e sensibilidade à topologia de
memória que tornam o agendamento ingênuo de compactações potencialmente prejudicial.
Compreender a interação entre a mecânica da LSM-tree e as características da PMem
é, portanto, central para a proposta deste trabalho.

\section{Memória Persistente (PMem)}
\label{sec:pmem}

A Memória Persistente (PMem), materializada em tecnologias como a Intel Optane DC
Persistent Memory, introduz uma nova camada na hierarquia de memória situada entre
a DRAM e os SSDs NVMe. Ao contrário da DRAM, que perde seus dados quando a energia
é removida, e ao contrário dos SSDs, que são acessíveis apenas por meio de
operações de bloco (4~KB), a PMem oferece simultaneamente \textbf{persistência de
dados} e \textbf{endereçamento a byte} (\textit{byte-addressability}), com latências
próximas à DRAM~\cite{van_renen_2020}. Essa posição singular na hierarquia tem
implicações diretas para o design de \textit{engines} de armazenamento: operações
que antes exigiam uma travessia completa pela pilha de I/O do sistema operacional
podem, em PMem, ser executadas diretamente pelo processador via instruções de carga
e armazenamento (\texttt{MOV}).

Trabalhos recentes exploraram especificamente as implicações da PMem para sistemas
baseados em LSM-tree. \citeauthor{revisting_design_lsm}~\cite{revisting_design_lsm}
revisitam o design de motores OLTP sobre PMem, identificando que a assimetria de
largura de banda é o principal fator limitante. \citeauthor{MatrixKV}~\cite{MatrixKV}
propõem um contêiner em NVM para reduzir \textit{write stalls}, e
\citeauthor{FlatLSM}~\cite{FlatLSM} apresentam uma variante de LSM-tree otimizada
para escrita em PMem. \citeauthor{SLM-DB}~\cite{SLM-DB} propõem uma abordagem de
nível único (\textit{single-level}) que aproveita o endereçamento a byte da PMem
para eliminar níveis intermediários. Esses trabalhos evidenciam que o uso eficiente
da PMem exige adaptações profundas nas políticas de gerenciamento da LSM-tree, o que motiva diretamente a proposta de escalonamento adaptativo deste trabalho.

\subsection{Modos de Operação}

A Intel Optane suporta dois modos principais de operação, com características
distintas para cada caso de uso~\cite{van_renen_2020}:

\begin{itemize}
    \item \textbf{Memory Mode:} A PMem funciona como extensão transparente da
    DRAM, gerenciada pelo controlador de memória. A DRAM atua como cache de
    última camada para a PMem. Não há persistência visível ao software. Este modo
    é útil para ampliar a capacidade de memória principal sem modificações no
    software.

    \item \textbf{App Direct Mode:} A PMem é exposta como um dispositivo separado,
    diretamente acessível pelo aplicativo via \textit{memory-mapped files} (usando
    bibliotecas como a \texttt{libpmem} do PMDK). Nesse modo, o software é
    responsável por garantir a durabilidade dos dados, gerenciando explicitamente
    os \textit{flushes} de linha de cache.
\end{itemize}

\subsection{Domínios de Persistência e Primitivas}

A garantia de durabilidade na PMem exige o gerenciamento explícito dos caches da
CPU. Os dados escritos por uma instrução \texttt{MOV} podem residir nos caches
voláteis (L1/L2/L3) indefinidamente. Para persistir, o dado deve alcançar o
domínio de persistência assíncrona (ADR) no controlador de memória, que garante
que os dados serão gravados na PMem mesmo em caso de falha de
energia~\cite{van_renen_2020}.

As instruções necessárias para este processo são:

\begin{itemize}
    \item \textbf{Cache Line Flush:} Instruções como \texttt{CLWB}
    (\textit{Cache Line Write Back}) escrevem a linha de cache suja de volta para
    a memória sem necessariamente invalidá-la no cache, o que é crucial para
    manter o desempenho de leituras subsequentes.

    \item \textbf{Fences (Barreiras):} A instrução \texttt{SFENCE} é obrigatória
    após os \textit{flushes} para garantir a ordem das operações e confirmar que
    os dados deixaram os buffers de escrita da CPU.
\end{itemize}

Cada \textit{persist-barrier} (par \texttt{CLWB} + \texttt{SFENCE}) tem um custo
de latência não desprezível~\cite{van_renen_2020}. Em cenários de alta
concorrência, múltiplas \textit{threads} emitindo \textit{persist-barriers}
simultaneamente podem saturar os buffers internos do controlador de PMem,
resultando em degradação severa de latência.

\subsection{Desempenho e Assimetria}

A avaliação experimental conduzida por~\citeauthor{van_renen_2020}~\cite{van_renen_2020}
revela comportamentos críticos para o escalonamento:

\begin{itemize}
    \item \textbf{Assimetria Leitura/Escrita:} A largura de banda de escrita na
    PMem é significativamente menor que a de leitura, podendo ser de 3 a 4
    vezes inferior em cargas aleatórias~\cite{van_renen_2020}. Isso significa que
    compactações agressivas, que são essencialmente cargas de escrita intensiva,
    consomem uma fração desproporcional da largura de banda
    disponível~\cite{revisting_design_lsm}.

    \item \textbf{Granularidade Interna:} Embora endereçável a byte, a PMem opera
    internamente com blocos maiores (256~bytes no caso da Intel Optane). Escritas
    muito pequenas (menores que 256~B) causam amplificação interna de escrita no
    dispositivo (\textit{media write amplification}), penalizando tanto o
    desempenho quanto a durabilidade~\cite{van_renen_2020}.

    \item \textbf{Sensibilidade NUMA:} A PMem é um dispositivo topologicamente
    vinculado a um soquete NUMA específico. Acessos realizados por \textit{threads}
    alocadas em soquetes remotos atravessam os enlaces \textit{inter-socket} (UPI),
    introduzindo latência adicional significativa.

    \item \textbf{Contenção sob Alta Concorrência:} Um excesso de \textit{threads}
    tentando persistir dados simultaneamente pode saturar os buffers do controlador
    de memória, degradando drasticamente a latência~\cite{van_renen_2020}. Essa
    característica fundamenta a necessidade de um escalonador que limite a
    concorrência de compactações em momentos de alta pressão de I/O, conforme
    investigado também em~\cite{MatrixKV,FlatLSM}.
\end{itemize}
